\chapter{Conclusion}
\label{ch:conclusion}

The problem this project set out to address is that guidance for visitors
in Nepal ends at the entrance: inside a heritage site or a campus there is
no structured navigation, no storytelling at the places where it matters,
no incentive to explore, and no way for an institution to verify that
anyone was actually there.

\projectname{} answers with a complete, working system rather than a
concept. A browser-based AR pathfinder guides visitors waypoint by
waypoint; a server-authoritative geofence with two-phase confirmation makes
every recorded visit evidence of physical presence; a gamification engine
--- points, milestones, badges, side quests, quizzes, and redeemable
rewards --- gives exploration a reason to continue and to repeat; and a
role-separated admin surface lets a non-technical institution configure
routes, media, quizzes, and geofence radii without touching code.

The build was validated where it counts: on the ground. Real users walked
the seeded Orchid College route during the hackathon, triggered geofenced
arrivals, watched story media anchored in the camera view, earned badges
and points exactly once each, and appeared on the leaderboard --- while the
admin dashboard accumulated engagement data as a side effect of ordinary
use. Because every site-specific element is configuration rather than
code, the same platform that guided a campus walk can be pointed at any
museum or heritage property with an admin account and a set of
coordinates. The prototype thus demonstrates the full loop it promised:
verified presence, engaging interpretation, and measurable exploration,
delivered as software alone.
